\chapter{Templates}\label{app:templates}

Copy these into your repository (\file{docs/templates/}) and use them as-is.
They are deliberately short: a template you actually fill in beats a perfect one
you skip.

\section*{Lab notebook entry --- \texttt{docs/notebook/YYYY-MM-DD.md}}
\begin{lstlisting}[style=plain]
# 2027-02-14 (week 5, Phase 1)

Goal today:      log-domain trackers -- does 8-bit storage meet RQ1?
Starting commit: 3f2a9c1

## Did
- Implemented fx_imcra.log_tracker (bias-encoded log2, 8 bits)
- Ran experiments/023-log-trackers on dev subset (200 utts)

## Results
- dPESQ = -0.03 [-0.04, -0.02] vs float; dSTOI = -0.002
- Noise PSD median error 0.31 dB

## Problems / surprises
- Edge bins k=0,1 biased high -- check smoothing zero-padding

## Next
- Fix edge bins; rerun; then joint run with 16-bit S, S~, lambda
\end{lstlisting}

\section*{Experiment record --- \texttt{experiments/NNN-name/notes.md}}
\begin{lstlisting}[style=plain]
# Experiment 023: log-domain minimum trackers

Question:    Can the tracker minima be stored as 8-bit log2 values?
Hypothesis:  Yes; |dPESQ| <= 0.05 on the dev set (part of H1)
Config:      config.yaml (this folder)     Commit: 3f2a9c1    Seed: 0
Data:        dev subset, 200 utterances (list in dev_ids.txt)
Command:     python -m golden.sweep --config experiments/023-log-trackers/config.yaml
Results:     results.csv (per utterance), summary.md (means + 95% CIs)
Conclusion:  Supported / not supported, because ...
Next step:   ...
\end{lstlisting}

\section*{Weekly status report (one page, for your supervisor)}
\begin{lstlisting}[style=plain]
Week 9 of 26 -- Phase 2 (RTL)        Plan status: on track / 1 week behind / ...

Done this week:     - recip_unit bit-exact (V2), min_tracker bit-exact (V3)
Planned, not done:  - gain_lut commit test (blocked: AXI write ordering bug)
Next week:          - fix AXI bug; start bin_pipeline stages 1-3
Risks (from register): R2 FFT complexity -- unchanged; R7 Vivado host -- resolved
Decision needed:    - LR or SVM head for Track A? (data: Phase 1 table 3)
\end{lstlisting}

\section*{Design review checklist (before a block is called done)}
\begin{checklist}
  \item Ports, widths and formats match \file{docs/spec.md}.
  \item Lint-clean with \code{-Wall}; no inferred latches; memories inferred as RAM.
  \item One reset style; all registers reset or deliberately left without reset.
  \item Rounding and saturation via the shared helpers only.
  \item Parameters read by the datapath are shadowed and latched at frame boundaries.
  \item A testbench exists, passes on two simulators, and appears in the verification plan.
  \item Resource numbers recorded in \file{docs/resources.md}.
\end{checklist}

\section*{Verification plan row}
\begin{center}\footnotesize
\begin{tabularx}{\textwidth}{@{}lLLLl@{}}
\toprule
ID & Feature & Test (file) & Pass criterion & Status \\
\midrule
V\textit{n} & & \file{tb/...} & & planned / passing / failing \\
\bottomrule
\end{tabularx}
\end{center}

\section*{Specification row}
\begin{center}\footnotesize
\begin{tabularx}{\textwidth}{@{}llllL@{}}
\toprule
Signal & Symbol & Format & Range (measured) & Evidence (experiment) \\
\midrule
 & & Q\textit{m}.\textit{F} & & \file{experiments/NNN-...} \\
\bottomrule
\end{tabularx}
\end{center}

\section*{Milestone report (one page)}
\begin{lstlisting}[style=plain]
Milestone M2 -- RTL bit-exact in simulation        Date:        Tag: m2-bit-exact

Exit criteria (from Chapter 13)          Evidence
- [x] all blocks lint-clean               CI run #212
- [x] ise_top bit-exact, 6 utterances     tb/top/test_utterances.py, log
- [ ] coverage gaps explained             2 open items (see verification_plan.md)

What went well / what was hard:
Deviations from plan (and why):
Changes to plan for next phase:
\end{lstlisting}
